\subsection{负阶 Sobolev 空间与对偶性}\label{sec:negative-sobolev-duality}

根据 Riesz 表示定理, Hilbert 空间与其对偶空间等距. 因而 $H^k(\Omega)$ 上每个线性泛函都具有 $l(v)=(u,v)_k$ 形式的表示. 然而, 内积 $(u,v)_k$ 与 $u$ 作为分布的作用并不一致. 事实上, 由于试验函数一般在 $H^k(\Omega)$ 中不稠密, 线性泛函未必是分布; 存在对所有试验函数都为零的非零线性泛函. 作如下定义.

\begin{Definition}
\label{def:negative-order-sobolev-dual-spaces}
用 $H^{-k}(\Omega)$ 表示 $H_0^k(\Omega)$ 上所有线性泛函的集合. 此外, 若 $M$ 为 $\mathbb{R}^m$ 或 $C^k$ 类紧流形, $k>s$, 则 $H^{-s}(M)$ 表示 $H^s(M)$ 的对偶空间.
\end{Definition}

由于 $\mathcal{D}(\Omega)$ 在 $H_0^k(\Omega)$ 中稠密, $H^{-k}(\Omega)$ 是一个分布空间. 正如下面的例子所示, 负阶 Sobolev 空间包含奇异分布.

\Needspace{4\baselineskip}
\begin{Example}
\label{ex:dirac-negative-sobolev-membership}
假设 $k>m/2$, $\Omega\subset\mathbb{R}^m$ 具有 $k$ 阶延拓性质且包含原点. 那么 Dirac delta 属于 $H^{-k}(\Omega)$. 为说明这一点, 注意 Sobolev 嵌入定理保证 $H^k(\Omega)$ (因而还有 $H_0^k(\Omega)$) 连续嵌入 $C_b(\overline\Omega)$. 这保证 delta 分布有良好定义. 它也是 $H_0^k(\Omega)$ 上的有界线性泛函, 因为对每个 $u\in H_0^k(\Omega)$ 有\[|(\delta,u)|:=|u(\bm{0})|\leq\|u\|_{\infty}\leq C\|u\|_{H^k(\Omega)}.\]
\end{Example}

\Needspace{4\baselineskip}
\begin{Example}
\label{ex:surface-charge-negative-sobolev}
设 $S$ 为 $\Omega\subset\mathbb{R}^3$ 内部的光滑有界曲面, $g:S\to R$ 属于 $L^2(S)$. (可将 $g$ 看作 $S$ 上的面电荷分布.) 那么由 $g$ 生成的分布属于 $H^{-1}(\Omega)$. 为说明这一点, 利用迹定理注意到对任意光滑函数 $\phi$ 有\[\begin{aligned}
|(g,\phi)|&=\left|\int_Sg(\bm{x})\phi(\bm{x})\,\mathrm{d}A(\bm{x})\right|\\
&\leq\|g\|_{L^2(S)}\|\phi\|_{L^2(S)}\\
&\leq\|g\|_{L^2(S)}\|\phi\|_{H^{1/2}(S)}\\
&\leq\|g\|_{L^2(S)}\|\phi\|_{H^1(\Omega)}
\end{aligned}\]因此, 面分布 $g$ 定义了 $H^1(\Omega)$ 中函数上的有界线性泛函.
\end{Example}

也可将负阶 Sobolev 空间中的函数刻画为正阶 Sobolev 空间中函数的导数. 设 $f\in H^{-k}(\Omega)$. 由 Riesz 表示定理, 存在唯一的 $u\in H_0^k(\Omega)$, 使对每个 $v\in H_0^k(\Omega)$ 有
\begin{equation}\label{eq:negative-sobolev-riesz-representation}
(f,v)=(u,v)_k
\end{equation}
$u$ 与 $f$ 有什么关系? 由 \eqref{eq:negative-sobolev-riesz-representation} 可知, 对任意试验函数 $\phi$,
\begin{equation}\label{eq:negative-sobolev-riesz-derivative-pairings}
(f,\phi)=\sum_{|\alpha|\leq k}(D^\alpha u,D^\alpha\phi)=\sum_{|\alpha|\leq k}(-1)^{|\alpha|}(D^{2\alpha}u,\phi),
\end{equation}
即
\begin{equation}\label{eq:negative-sobolev-elliptic-derivative-representation}
f=\sum_{|\alpha|\leq k}(-1)^{|\alpha|}D^{2\alpha}u.
\end{equation}


因此, 对任意给定的 $f\in H^{-k}(\Omega)$, 存在唯一的 $u\in H_0^k(\Omega)$ 满足偏微分方程 \eqref{eq:negative-sobolev-elliptic-derivative-representation}. 回顾条件 $u\in H_0^k(\Omega)$ 可解释为边界条件: 在 $\partial\Omega$ 上有 $u=\partial u/\partial n=\cdots=\partial^{k-1}u/\partial n^{k-1}=0$ (Theorem~\ref{thm:zero-boundary-sobolev-characterization}). 与上述类似的考察构成椭圆边值问题现代存在性理论的出发点; 将在 Chapter~\ref{chap:linear-elliptic-equations} 回到这一问题.

最后给出一个关于负阶 Sobolev 空间中分布求导的简单断言.

\begin{Lemma}
\label{lem:integer-order-sobolev-derivative}
设 $u\in H^k(\Omega)$, $k\in\mathbb{Z}$. 那么 $\partial u/\partial x_i\in H^{k-1}(\Omega)$.
\end{Lemma}

证明由定义显然得到. Lemma~\ref{lem:integer-order-sobolev-derivative} 有一个逆命题.

\begin{Lemma}
\label{lem:negative-sobolev-l-two-derivative-representation}
设 $f\in H^{-k}(\Omega)$, $k\in\mathbb{N}$. 那么存在函数 $g_\alpha\in L^2(\Omega)$, 使 $f=\sum_{|\alpha|\leq k}D^\alpha g_\alpha$.
\end{Lemma}

证明时, 只需在 \eqref{eq:negative-sobolev-elliptic-derivative-representation} 中令 $g_\alpha=(-1)^{|\alpha|}D^\alpha u$.
